\title{FlashAttention-3: Fast and Accurate Attention with Asynchrony and Low-precision}

\begin{abstract}
Attention serves as a fundamental mechanism within the Transformer framework, yet it constitutes a critical bottleneck for scaling large language models and supporting long-context scenarios. \textsc{FlashAttention} introduced a GPU-optimized strategy that reduces overhead by minimizing memory access operations. Nevertheless, this method does not harness the advanced features available in state-of-the-art hardware, resulting in \textsc{FlashAttention-2} utilizing only 35\% of H100 GPU capacity. In response, we propose three principal innovations to accelerate attention on Hopper GPUs: leveraging the asynchronous operation of Tensor Cores and TMA for (1) maximizing concurrency between computation and data transfers using warp-specialization, (2) interweaving the execution of block-wise matrix multiplication and softmax, and (3) implementing block-wise quantization and processing methods adapted to FP8 low-precision, made possible by hardware enhancements. Experimental results indicate our approach, \textsc{FlashAttention-3}, achieves a 1.5-2.0$\times$ throughput improvement on H100 GPUs, with FP16 performance scaling up to 740 TFLOPs/s (corresponding to 75\% utilization), and FP8 performance nearing 1.2 PFLOPs/s. Furthermore, evaluations reveal that FP8 \textsc{FlashAttention-3} produces 2.6$\times$ less numerical error compared to a baseline FP8 attention mechanism.
\end{abstract}

\section{Introduction}
\label{sec:intro}

Within the Transformer architecture~\citep{vaswani2017attention}, the attention module represents the central computational bottleneck, owing to the quadratic time complexity of calculating self-attention scores between queries and keys as a function of sequence length. Enabling attention over longer sequences could facilitate important advances, such as improved modeling and reasoning across extensive textual corpora~\citep{guo2021longt5,shaham2022scrolls,peng2023yarn} or processing numerous files in extensive code repositories~\citep{roziere2023code, li2023starcoder}, as well as expanding to additional data modalities like high-resolution images~\citep{chen2022scaling}, audio signals~\citep{gulati2020conformer}, and video streams~\citep{ho2022video}. Moreover, extended attention empowers new practical applications, for instance handling user sessions with substantial history~\citep{sun2019bert4rec} or supporting agent-driven workflows with extended planning horizons~\citep{yao2022react}. As a result, substantial effort has focused on accelerating attention computations for long sequences, including approximate techniques~\citep{katharopoulos2020transformers,choromanski2020rethinking,
tay2020efficient}, software-level enhancements (\citep{rabe2021self,dao2022flashattention,kwon2023efficient}), and the development of fundamentally different architectures~\citep{peng2023rwkv,sun2023retentive,gu2023mamba}.

This study extends the advancements made by \citet{dao2022flashattention}, who developed exact-attention algorithms by incorporating detailed awareness of the GPU's execution architecture and hardware features within their algorithmic frameworks. In \citep{dao2022flashattention}, Dao et al. proposed \textsc{FlashAttention}, an innovative approach employing tiling for the parallelization of attention mechanisms, which removes redundant accesses to slower global memory by merging all attention-related operations inside a single GPU kernel. 

Building upon this, \citet{dao2023flashattention2} redesigned the method as \textsc{FlashAttention-2}, enabling parallelism along the sequence length dimension and restructuring the forward pass's inner loop to operate over blocks of the key and value matrices. This modification led to better work distribution and occupancy rates on GPUs. Despite these improvements, our analysis reveals that \textsc{FlashAttention-2} does not fully exploit newer GPUs, such as the Hopper H100, in comparison to highly-tuned matrix-multiplication (GEMM) kernels—recording only 35\% utilization versus 80-90\%. This suboptimal performance is partly due to choices at the implementation level, like relying on Ampere-specific instructions rather than those tailored for Hopper when employing Tensor Cores. Recent efforts including ThunkerKitten~\citep{spector2024thunder} and cuDNN 9~\citep{cudnn9} illustrate that integrating Hopper-specific instructions with tile-based methodologies not only accelerates attention computations but also streamlines their implementation.

At a foundational level, the algorithm underlying \textsc{FlashAttention-2} operates within a conventional synchronous framework and does not intentionally incorporate asynchrony or low-precision strategies into its architecture. Asynchronous execution arises from the specialization of hardware components, designed to accelerate key computations in machine learning tasks: for instance, dedicated modules such as Tensor Cores facilitate matrix multiplications, while the Tensor Memory Accelerator (TMA) handles memory transactions, working independently from general-purpose CUDA cores that manage various computational operations. The use of reduced numerical precision—examples include FP8 in Hopper and FP4 in Blackwell, building on earlier deployments of FP16 (Pascal, 2017) and BF16 (Ampere, 2020)—is an established approach for doubling or quadrupling throughput without increasing power consumption or die size. A discussion of Hopper’s advancements in these areas can be found in \cref{subsec:hardware}.
A central engineering problem is to adapt \textsc{FlashAttention-2} so it can exploit these advanced hardware capabilities: leveraging asynchrony entails interleaving matmul and softmax computations, even in scenarios where one operation is contingent upon the completion of another; meanwhile, supporting low-precision demands diligent techniques for preserving accuracy, notably in the presence of outlier features within LLMs~\citep{dettmers2208llm, sun2024massive}.

To address this objective, we introduce \textsc{FlashAttention-3}, integrating and advancing three novel concepts to boost efficiency on modern GPU platforms:\footnote{Our empirical findings focus primarily on NVIDIA's Hopper architecture. Nonetheless, the underlying algorithm remains applicable to any GPU environment supporting adequate asynchronous operations and low-precision functionality.}

\begin{enumerate}

\item \textbf{Producer-Consumer asynchrony:} Our approach implements a warp-specialized software pipeline that leverages the asynchronous behavior of both data transfer and Tensor Core computations. By assigning data producers and consumers to distinct warps, we enhance the algorithm's capability to mask latencies associated with memory access and instruction issuing.
\item \textbf{Hiding softmax under asynchronous block-wise GEMMs:} To maximize throughput, we schedule the relatively slower non-GEMM steps required by softmax—such as floating-point multiply-add and exponential computations—concurrently with asynchronous WGMMA GEMM instructions. This involves a restructuring of the \textsc{FlashAttention-2} method to eliminate some of the intrinsic sequential dependencies between softmax operations and GEMMs. Specifically, in the two-stage variant, softmax is performed on a single block of the score matrix while WGMMA operates asynchronously to prepare the subsequent block.
\item \textbf{Hardware-accelerated low-precision GEMM:} We revise the forward pass algorithm to utilize the FP8 Tensor Cores for GEMM computations, thereby almost doubling the attained TFLOPs/s. Accommodating this requires addressing the disparate memory layouts expected by WGMMA for blocks of FP32 accumulators and FP8 operand matrices. The adoption of block quantization and incoherent processing strategies helps alleviate accuracy degradation inherent to transitioning to FP8 precision.
\end{enumerate}

To assess the empirical performance of our approach, we conducted extensive benchmarking of \textsc{FlashAttention-3} on the H100 SXM5 GPU across diverse settings. Our results demonstrate that (1) the FP16 variant yields a 1.5–2.0$\times$ acceleration over \textsc{FlashAttention-2} during the forward pass (with performance peaking at 740 TFLOPs/s) and a 1.5–1.75$\times$ improvement in the backward pass, (2) FP8 reaches nearly 1.2 PFLOPs/s, and (3) for scenarios involving lengthy sequences, FP16 offers superior throughput while FP8 remains competitive\footnote{Specifically, \textsc{FlashAttention-3} FP8 surpasses others for a head dimension of 64; for head dimensions of 128 and 256, it matches alternatives in cases without causal masking but falls behind with causal masking.} when compared to NVIDIA's cuDNN attention implementation. Furthermore, we confirm that FP16 \textsc{FlashAttention-3} exhibits identical numerical error to \textsc{FlashAttention-2}, outperforming conventional attention computations by preserving intermediate values (e.g., softmax scaling) in FP32 precision. Notably, FP8 \textsc{FlashAttention-3}, which employs block quantization alongside incoherent execution, achieves 2.6$\times$ higher accuracy compared to standard attention that uses per-tensor quantization, particularly in settings involving outlier features.

We have released \textsc{FlashAttention-3} as open source under a permissive license\footnote{\textsc{FlashAttention-3}
  is accessible at \texttt{https://github.com/Dao-AILab/flash-attention}}, and aim to incorporate it into PyTorch and Hugging Face libraries so that it can be leveraged by as many researchers and developers as possible.

\section{Background: Multi-Head Attention and GPU Characteristics}
\label{sec:background}

\subsection{Multi-Head Attention}
\label{subsec:multi_head_attn}

Let $\mathbf{Q}, \mathbf{K}, \mathbf{V} \in \mathbb{R}^{N \times d}$ be the query, key and value input sequences associated to a single head, where $N$ is the sequence length and $d$ is the head dimension. Then the attention output $\mathbf{O}$ is computed as:

\begin{equation*}
  \mathbf{S} = \alpha \mathbf{Q} \mathbf{K}^\top \in \mathbb{R}^{N \times N}, \quad \mathbf{P} = \mathrm{softmax}(\mathbf{S}) \in \mathbb{R}^{N \times N}, \quad \mathbf{O} = \mathbf{P}\mathbf{V} \in \mathbb{R}^{N \times d},
\end{equation*}

The $\mathrm{softmax}$ function is performed on each row, and usually, the scaling coefficient $\alpha$ is chosen as $1/\sqrt{d}$. To mitigate numerical issues when dealing with exponentiation, it is standard practice to subtract $\mathrm{rowmax}(\mathbf{S})$ from $\mathbf{S}$. In the context of multi-head attention (MHA), distinct query, key, and value projection matrices are assigned to each attention head. This allows the calculations to be distributed over several heads and batch entries concurrently, resulting in the complete output tensor.

Now let $\phi$ be a scalar loss function and let $\mathbf{d}(-) = \partial \phi / \partial (-)$ be notation for the gradient. Given the output gradient $\mathbf{dO} \in \mathbb{R}^{N \times d}$, we compute $\mathbf{dQ}$, $\mathbf{dK}$, and $\mathbf{dV}$ according to the chain rule as follows:

\begin{align*}
  \mathbf{dV} &= \mathbf{P}^\top \mathbf{dO} \in \mathbb{R}^{N \times d} \\
  \mathbf{dP} &= \mathbf{dO} \mathbf{V}^\top \in \mathbb{R}^{N \times N} \\
  \mathbf{dS} &= \mathrm{dsoftmax} (\mathbf{dP}) \in \mathbb{R}^{N \times N} \\
  \mathbf{dQ} &= \alpha \mathbf{dS} \mathbf{K} \in \mathbb{R}^{N \times d} \\
  \mathbf{dK} &= \alpha \mathbf{dS}^\top \mathbf{Q} \in \mathbb{R}^{N \times d},
\end{align*}

In this context, $\mathbf{d}s = (\mathrm{diag}(p) - p p^\top)\mathbf{d}p$ is valid when $p = \mathrm{softmax}(s)$ with $s$ treated as a vector. We denote the application of this formula to each row as $\mathrm{dsoftmax}(\mathbf{dP})$. Ultimately, the backward pass for MHA can be executed concurrently across all heads and batches using this approach.

\subsection{GPU hardware characteristics and execution model}
\label{subsec:hardware}

We outline key elements of the GPU execution model that pertain to \textsc{FlashAttention-3}, emphasizing the NVIDIA Hopper architecture as a specific example to ground our discussion.

\paragraph{Memory hierarchy:} The architecture of the GPU’s memory is structured as a hierarchy of data locations, where larger capacities come with lower bandwidth, as detailed in \cref{tab:gpu-hierarchy}\footnote{Shared memory bandwidth is given as 128 bytes per clock cycle per SM in \citet{luo2024benchmarking}. We calculate total bandwidth by multiplying this figure by 132 SMs and a boost clock frequency of 1830 MHz.}. At the highest level, global memory (GMEM), sometimes referred to as HBM, serves as off-chip DRAM that can be accessed by all streaming multiprocessors (SMs). Data transferred from GMEM is automatically cached within the on-chip L2 cache. Furthermore, each SM is equipped with a compact, on-chip, heavily banked shared memory (SMEM), which is explicitly managed by the programmer. The lowest layer of the hierarchy is formed by the register file located within each SM.

\paragraph{Thread hierarchy:} In the GPU programming paradigm, execution units known as threads are structured into various logical groupings. Starting from the smallest, the hierarchy includes threads, followed by warps—which consist of 32 threads each—then warpgroups, made up of 4 sequential warps, and threadblocks (also referred to as cooperative thread arrays or CTAs). On newer architectures like Hopper, threadblock clusters exist above threadblocks, and at the highest level, all these groups together form grids.

The two hierarchies are tightly connected. Threads belonging to a single CTA are simultaneously scheduled on a common SM, while CTAs assigned to a given cluster share scheduling on a single GPC. All threads in a CTA can directly access SMEM, but each thread has its own RMEM, limited to a maximum of 256 registers.

\paragraph{Asynchrony and warp-specialization:}

GPUs are designed as throughput-oriented processors, leveraging parallelism and asynchronous operations to mitigate latency in memory access and computation. In Hopper, asynchronous data transfer between GMEM and SMEM is facilitated by the Tensor Memory Accelerator (TMA), which is implemented as a specialized hardware component \cite[\S7.29]{cuda}. Additionally, Hopper distinguishes itself from earlier architectures like Ampere by offering a Tensor Core capable of operating asynchronously. This Tensor Core is accessible through the warpgroup-wide WGMMA instruction \cite[\S9.7.14]{ptx}, and uniquely, it can retrieve its input data directly from shared memory.

With hardware-enabled asynchrony, kernels can be designed for warp specialization, assigning each warp within a CTA to distinct producer or consumer functions—restricted solely to data transfer or computation, respectively. This structural separation facilitates more effective generation of instruction schedules by the compiler \citep{warp-specialization-2011}. Furthermore, Hopper introduces the capability for registers to be dynamically redistributed among warpgroups through the use of \verb|setmaxnreg| \cite[\S9.7.17.1]{ptx}, enabling warps engaged in MMAs to access a greater proportion of RMEM compared to those limited to executing TMA, which require only a single participating thread.

\paragraph{Low-precision number formats:}
\label{sec:low-precision-gpu}
Contemporary GPUs incorporate dedicated hardware components designed to enhance the efficiency of low-precision arithmetic operations. As an illustration, the WGMMA instruction allows execution on FP8 Tensor Cores within the Hopper architecture, providing double the throughput per streaming multiprocessor (SM) relative to FP16 or BF16 precision modes.

Accurately employing FP8 WGMMA requires familiarity with the operand layout restrictions. Consider a GEMM operation computing $A \times B^{\top}$, with $A$ of size $M\times K$ and $B$ of size $N\times K$. The $A$ or $B$ operand is referred to as \emph{mn-major} when it has memory contiguity along the $M$ or $N$ axis, whereas it is called \emph{k-major} if contiguous along the $K$ axis. In the case of FP16 WGMMA, both mn-major and k-major representations for input operands in SMEM are permitted. Conversely, FP8 WGMMA only supports the k-major format. Additionally, scenarios like attention, where sequential GEMM operations are fused within the same kernel, are complicated by mismatched layouts between FP32 accumulators and FP8 operands, making it difficult to utilize dependent FP8 WGMMAs.

Regarding attention mechanisms, these layout constraints necessitate specific adjustments in the construction of an FP8 algorithm, as elaborated in \cref{sec:algofp8}.

\subsection{Standard Attention and Flash Attention} In alignment with \citet{dao2022flashattention}, we use the term \textbf{standard attention} to refer to the approach on GPU where the intermediate matrices $\mathbf{S}$ and $\mathbf{P}$ are explicitly stored in HBM. The essential innovation of \textsc{FlashAttention} involves exploiting a localized variant of softmax reduction, thereby eliminating costly intermediate data transfers and integrating the attention computation within a single kernel operation. This localized softmax process is performed in lines \ref{code-ws:softmax_start}-\ref{code-ws:softmax_end} of the consumer mainloop from \cref{alg:flash3_wgmma_ws_only}, in conjunction with rescaling of the $\mathbf{O}$ matrix blocks. The justification that this approach correctly evaluates $\mathbf{O}$ is provided in \cite[\S 2.3.1]{dao2023flashattention2}.

\section{FlashAttention-3: Algorithm}
\label{sec:algo}

This section introduces the \textsc{FlashAttention-3} algorithm. To streamline the exposition, we concentrate on the forward pass; details on the backward pass are provided in~\cref{sec:algo_ws_bwd}. Our discussion begins with an overview of integrating warp-specialization, leveraging a circular SMEM buffer, into the foundational procedure of \textsc{FlashAttention-2}. Next, we demonstrate how the asynchrony inherent in WGMMA can be harnessed to establish a two-stage pipeline where GEMM and softmax computations are overlapped. Lastly, we outline the adjustments necessary for FP8, encompassing layout compliance as well as enhancements to accuracy using block quantization and incoherent processing techniques.

\subsection{Producer-Consumer asynchrony through warp-specialization and pingpong scheduling}
\label{sec:algo_ws}

\paragraph{Warp-specialization}
Similar to \textsc{FlashAttention-2}, the forward computation in \textsc{FlashAttention-3} exhibits a high degree of parallelism with respect to batch size, head count, and the length of the query sequence. Therefore, analyzing the algorithm at the CTA level is sufficient, where the algorithm processes a tile $\mathbf{Q}_i$ from the query matrix and generates the corresponding tile $\mathbf{O}_i$ in the output. For clarity, we first describe the warp-specialization strategy using a circular SMEM buffer, intentionally omitting the overlap of GEMM and softmax operations for now. Denote $d$ as the dimension per head, $N$ as the sequence length, and choose a query block size $B_r$ such that the query matrix $\mathbf{Q}$ is partitioned into $T_r = \lceil \frac{N}{B_r} \rceil$ blocks, namely, $\mathbf{Q}_1, ..., \mathbf{Q}_{T_r}$.

\begin{algorithm}[H]
    \caption{\small\label{alg:flash3_wgmma_ws_only}\textsc{FlashAttention-3} forward pass \textbf{excluding} intra-consumer overlapping -- CTA perspective}
    \begin{algorithmic}[1]
\REQUIRE Input matrices $\mathbf{Q}_i \in \mathbb{R}^{B_r \times d}$, $\mathbf{K}, \mathbf{V} \in \mathbb{R}^{N \times d}$ residing in HBM; specified key block size $B_c$, with $T_c = \lceil \frac{N}{B_c} \rceil$.
\STATE Create pipeline structure for barrier synchronization using a $s$-stage circular SMEM buffer.
\IF {current warpgroup is producer}
\STATE Release a set quantity of registers in advance.
\STATE Transfer $\mathbf{Q}_i$ from HBM into shared memory.
\STATE After the transfer, update status to indicate $\mathbf{Q}_i$ is ready for consumers.
\FOR{$0 \le j < T_c$}
    \STATE Await buffer consumption at the $(j\,\%\,s)$th pipeline stage.
    \STATE Transfer $\mathbf{K}_j$ and $\mathbf{V}_j$ from HBM into shared memory at buffer position $(j\,\%\,s)$.
    \STATE Signal consumers that $\mathbf{K}_j$ and $\mathbf{V}_j$ are now available.
\ENDFOR
\ELSE \STATE Acquire the released register pool, scaled by consumer warp count.
\STATE Initialize, on chip, the values $\mathbf{O}_i = (0) \in \mathbb{R}^{B_r \times d}$, and $\ell_i, m_i = (0), (-\infty) \in \mathbb{R}^{B_r}$.
\STATE Wait until $\mathbf{Q}_i$ becomes accessible in shared memory.
\FOR{$0 \le j < T_c$}
\STATE Block until $\mathbf{K}_j$ is present in shared memory.
\STATE Calculate $\mathbf{S}_i^{(j)} = \mathbf{Q}_i \mathbf{K}_j^T$ using SS-GEMM. Commit results and pause for completion. \label{alg:ws_only_gemm1}
\STATE Cache $m_i^{\mathrm{old}} = m_i$ and update $m_i = \mathrm{max}(m_i^{\mathrm{old}}, \mathrm{rowmax}(\mathbf{S}_i^{(j)}))$. \label{code-ws:softmax_start}
\STATE Calculate $\widetilde{\mathbf{P}}_i^{(j)} = \mathrm{exp}(\mathbf{S}_i^{(j)} - m_i)$ and revise $\ell_i = \mathrm{exp}(m_i^{\mathrm{old}} - m_i) \ell_i + \mathrm{rowsum}(\widetilde{\mathbf{P}}_i^{(j)})$. \label{code-ws:softmax_end}
\STATE Await until $\mathbf{V}_j$ is ready in shared memory.
\STATE Update $\mathbf{O}_i = \mathrm{diag}(\mathrm{exp}(m_i^{\mathrm{old}} - m_i))^{-1} \mathbf{O}_i + \widetilde{\mathbf{P}}_i^{(j)} \mathbf{V}_j$ with RS-GEMM. Commit and synchronize. \label{alg:ws_only_gemm2}
\STATE Mark the $(j\,\%\,s)$th buffer slot as free for the producer to reuse.
\ENDFOR
\STATE Apply the normalization: $\mathbf{O}_i = \mathrm{diag}(\ell_i)^{-1} \mathbf{O}_i$; compute $L_i = m_i + \log(\ell_i)$.
\STATE Store the outputs $\mathbf{O}_i$ and $L_i$ into HBM as block $i$ within $\mathbf{O}$ and $L$.
\ENDIF
\end{algorithmic}
\end{algorithm}

In our deployment of \cref{alg:flash3_wgmma_ws_only} on Hopper, we employ \verb|setmaxnreg| to manage (de)allocation tasks, utilize TMA for fetching $\mathbf{Q}_i$ along with $\{ \mathbf{K}_j, \mathbf{V}_j \}_{0 \leq j < T_c}$, and leverage WGMMA within the consumer mainloop for performing the GEMMs. The SS or RS prefix denotes whether the initial operand is retrieved from shared memory or the register file, respectively. When analyzing the operational sequence of \cref{alg:flash3_wgmma_ws_only}, it is important to recognize that TMA load instructions can be dispatched asynchronously and do not block pending completion of prior load requests. Additionally, during the producer mainloop, waits are absent in the first $s$ passes, allowing the buffer to accumulate data.

\paragraph{Pingpong scheduling} The asynchronous execution of WGMMA and TMA, coupled with warp-specialization, enables the possibility to interleave the softmax computation of one warpgroup with the GEMM operations in another. To illustrate the rationale behind this, consider that on current hardware accelerators, non-matmul tasks exhibit significantly less throughput compared to matmul tasks. For instance, the H100 SXM5 GPU offers 989 TFLOPS for FP16 matmul operations, whereas special functions such as exponential\footnote{According to the CUDA programming guide, each streaming multiprocessor (SM) can carry out 16 special function operations per clock cycle. By calculating $16 \times 132$ SMs $\times$ 1830 MHz, the resulting throughput is 3.9 TFLOPS for special functions.} (which are essential for softmax) reach only 3.9 TFLOPS. During the forward pass for attention in FP16 with a head dimension of 128, the required matmul FLOPS are 512 times greater than those needed for exponential operations, and since exponential functions deliver 256 times less throughput, they might occupy roughly 50\% of the runtime relative to matmul. This performance disparity becomes even more pronounced in FP8, where matmul throughput is doubled, but the exponential throughput remains unchanged.

Because the computation of the exponential is handled by a distinct hardware component (the multi-function unit), it is most efficient to initiate the exponential operation concurrently with the matmul executed by the Tensor Cores. To achieve this, we employ synchronization barriers (\texttt{bar.sync} instructions) to enforce an execution order: the GEMMs (GEMM1, representing $\mathbf{P} \mathbf{V}$ in one iteration, and GEMM0, corresponding to $\mathbf{Q} \mathbf{K}^\top$ in the following iteration) for warpgroup 1 are prioritized before those of warpgroup 2. Consequently, while warpgroup 2 computes its GEMMs, warpgroup 1 is simultaneously processing the softmax operation. The process then alternates, such that warpgroup 2 computes the softmax as warpgroup 1 returns to its GEMMs—implementing what is termed a ``pingpong'' schedule. This concept is visualized in~\cref{fig:pingpong_scheduling}. 

Although this pingpong scheduling does not occur in as strictly alternating a manner as the illustration might suggest, we have generally observed a notable enhancement in performance using this approach (for instance, performance increased from 570 TFLOPS to between 620 and 640 TFLOPS for FP16 forward computation with a head dimension of 128 and sequence length of 8192).

\paragraph{Attention variants} In the case of multi-query attention~\citep{shazeer2019fast} and grouped query attention~\citep{ainslie2023gqa}, our implementation adopts the strategy used in \textsc{FlashAttention-2}. Specifically, we modify the tensor indexing scheme so that $\mathbf{K}$ and $\mathbf{V}$ are not redundantly stored in HBM.

\subsection{Intra-warpgroup overlapping GEMMs and softmax}

It is possible to interleave certain instructions from the softmax operation with those from the GEMMs within a single warpgroup. Here, we outline a specific method to achieve this concurrency.

In the attention algorithm, operations within the inner loop (main loop) have sequential dependencies that impede parallelization within a single iteration. For example, (local) softmax (lines \ref{code-ws:softmax_start} to \ref{code-ws:softmax_end}) relies on the output $\mathbf{S}_i^{(j)}$ of the first GEMM, while the second GEMM takes its result $\widetilde{\mathbf{P}}_i^{(j)}$ as an operand. Indeed, the wait statements in lines \ref{alg:ws_only_gemm1} and \ref{alg:ws_only_gemm2} of \cref{alg:flash3_wgmma_ws_only} serialize the execution of softmax and GEMMs. However, we can break these dependencies by pipelining across iterations through additional buffers in registers. Pursuing this idea, we propose the following two-stage\footnote{Note that the number of stages of the overlapping scheme is bounded by, but need not equal, the number $s$ of stages in the circular SMEM buffer.} GEMM-softmax pipelining algorithm:

\begin{algorithm}[H]
    \caption{\small\label{alg:flash3_wgmma}\textsc{FlashAttention-3} consumer warpgroup forward pass}
    \begin{algorithmic}[1]
      \REQUIRE Matrices $\mathbf{Q}_i \in \mathbb{R}^{B_r \times d}$, $\mathbf{K}, \mathbf{V} \in \mathbb{R}^{N \times d}$ are stored in HBM; a key block size $B_c$ is given such that $T_c = \lceil \frac{N}{B_c} \rceil$.
      \STATE Adjust the allocation of a fixed number of registers depending on the total number of consumer warps.
      \STATE Initialize $\mathbf{O}_i = (0) \in \mathbb{R}^{B_r \times d}$, and set $\ell_i, m_i = (0), (-\infty) \in \mathbb{R}^{B_r}$ within on-chip memory.
      \STATE Await loading of $\mathbf{Q}_i$ and $\mathbf{K}_0$ into shared memory.
      \STATE Execute the matrix product $\mathbf{S}_{\mathrm{cur}} = \mathbf{Q}_i \mathbf{K}_0^T$ via WGMMA, then commit and synchronize.
      \STATE Free the first buffer stage associated with $\mathbf{K}$.
      \STATE Use $\mathbf{S}_{\mathrm{cur}}$ to compute $m_{i}$, $\tilde{\mathbf{P}}_{\mathrm{cur}}$, and $\ell_{i}$; subsequently, rescale $\mathbf{O}_i$.
      \FOR{$1 \le j < T_c - 1$}
        \STATE Await transfer of $\mathbf{K}_j$ into shared memory.
        \label{code:mainloop_start}
        \STATE Calculate $\mathbf{S}_{\mathrm{next}} = \mathbf{Q}_i \mathbf{K}_{j}^T$ using WGMMA, commit operation but continue without synchronization. \label{code:first_wgmma}
        \STATE Await availability of $\mathbf{V}_{j-1}$ in shared memory.
        \STATE Apply WGMMA to compute and accumulate: $\mathbf{O}_{i} = \mathbf{O}_{i} + \tilde{\mathbf{P}}_{\mathrm{cur}} \mathbf{V}_{j-1}$, commit and continue without synchronization. \label{code:second_wgmma}
        \STATE Wait for completion of WGMMA computing $\mathbf{Q}_i \mathbf{K}_{j}^T$.
        \STATE From $\mathbf{S}_{\mathrm{next}}$, determine $m_{i}$, $\tilde{\mathbf{P}}_{\mathrm{next}}$, and $\ell_{i}$. \label{code:softmax}
        \STATE After WGMMA for $\tilde{\mathbf{P}}_{\mathrm{cur}} \mathbf{V}_{j-1}$ finishes, apply rescaling to $\mathbf{O}_i$.
        \STATE Release buffer stage $(j\,\%\,s)$ for $\mathbf{K}$ and buffer stage $(j-1\,\%\,s)$ for $\mathbf{V}$, respectively.
        \STATE Assign $\mathbf{S}_{\mathrm{cur}} \leftarrow \mathbf{S}_{\mathrm{next}}$.
        \label{code:mainloop_end}
      \ENDFOR \STATE Await loading of $\mathbf{V}_{T_c - 1}$ in shared memory.
      \STATE Perform
      $\mathbf{O}_{i} = \mathbf{O}_{i} + \tilde{\mathbf{P}}_{\mathrm{last}} \mathbf{V}_{T_c - 1}$ via WGMMA, then commit and synchronize.

\STATE Epilogue: Adjust $\mathbf{O}_{i}$ using $m_{i}$. Derive $L_{i}$ utilizing $m_{i}$ and $\ell_{i}$. 
Store $\mathbf{O}_{i}$ and $L_{i}$ in HBM, corresponding to the $i$-th segment of $\mathbf{O}$ and $L$. 
\end{algorithmic}
\end{algorithm}

\cref{alg:flash3_wgmma} is introduced as a substitute for the consumer pathway in \cref{alg:flash3_wgmma_ws_only}, thereby constructing the full FP16 version of the \textsc{FlashAttention-3} algorithm. Generally, we employ WGMMA to refer to asynchronous GEMM operations. In the mainloop, spanning lines \ref{code:mainloop_start} through \ref{code:mainloop_end}, the execution of the second WGMMA instruction during iteration $j$ (identified in line \ref{code:second_wgmma}) is executed concurrently with the softmax computations associated with iteration $j+1$ (shown in line \ref{code:softmax}).

Although the pipelined architecture depicted previously can theoretically enhance performance, practical constraints must also be acknowledged:

\paragraph{Compiler reordering} The provided pseudocode assumes a sequentially ideal execution. However, in practice, the NVCC compiler frequently modifies instruction ordering to achieve its own optimization objectives. Such rearrangements can interrupt the intended alternation between WGMMA and non-WGMMA instructions, potentially causing undesired interactions or limiting the anticipated performance improvement. Examination of the resulting SASS code confirms that the compiler indeed produces overlapped instruction sequences as anticipated (refer to Section~\ref{sec:2-stage-sass}).

\paragraph{Register pressure} Sustaining peak efficiency depends on minimizing register spilling. The 2-stage pipeline structure, however, imposes additional register requirements for managing intermediate states and facilitating communication across pipeline stages. For instance, it necessitates reserving space for an extra $\mathbf{S}_{\mathrm{next}}$, thereby increasing register usage per threadblock by $B_r \times B_c \times \text{sizeof}(\text{float})$. Such an elevated register footprint may constrain the adoption of larger block sizes—which themselves intensify register consumption. Profiling-based compromises are generally required to select a balanced configuration.

\paragraph{3-stage pipelining} Building on the principles of the 2-stage approach, we introduce a 3-stage pipeline that enables further overlap, now aligning the second WGMMA operation with the softmax computation. This extension promises greater utilization of Tensor Cores but intensifies register requirements, as managing another pipeline stage demands even more registers. This, in turn, complicates the balance between maximizing tile dimensions and the benefits of deeper pipelines. Comprehensive details and empirical results regarding the 3-stage method are presented in ~\cref{sec:3-stage}.

\subsection{Low-precision with FP8}
\label{sec:algofp8}

\textbf{Efficiency: layout transformations.} Executing the forward pass for \textsc{FlashAttention-3} using FP8 precision introduces unique difficulties related to layout compatibility that are not present with FP16.

To begin, it is important to observe that input tensors $\mathbf{Q}$, $\mathbf{K}$, and $\mathbf{V}$ are generally stored with contiguity along the head dimension. However, for the second GEMM operation under FP8 WGMMA with a k-major requirement, $\mathbf{V}$—specifically, the portions loaded into SMEM—must instead be contiguous along the sequence length dimension. Given that the TMA load operation preserves the original contiguous dimension, addressing this issue necessitates one of two approaches: (1) transposing $\mathbf{V}$ in GMEM as a preparatory measure, or (2) performing an in-kernel transpose of $\mathbf{V}$ tiles after they are brought into SMEM. If opting for the first approach, implementation can proceed by either (1a) incorporating the transpose into the epilogue of a prior stage such as rotary embedding, or (1b) invoking an independent preprocessing transpose kernel\footnote{With proper optimization, the transpose kernel can operate at speeds approaching the device's memory bandwidth~\citep{colfax_cutlass_transpose_2024}.}, thereby swapping the strides between sequence length and head dimensions. Nonetheless, (1a) proves challenging to adapt within typical library infrastructures, and (1b) incurs excessive overhead in memory-constrained settings like inference.

In our implementation of FP8 \textsc{FlashAttention-3}, we select approach (2). For performing the in-kernel transpose, we utilize the LDSM (\verb|ldmatrix|) and STSM (\verb|stmatrix|) instructions, which allow a warp of threads to cooperatively transfer data from SMEM to RMEM and vice versa in increments of 128 bytes.\footnote{The PTX documentation specifies that LDSM/STSM operate on $8 \times 8$ matrices with 16-bit elements \cite[\S 9.7.13.4.15-16]{ptx}. However, by packing two 8-bit values together, we can employ these instructions in FP8 settings. Still, the transposed variants of LDSM/STSM do not separate the packed 8-bit elements, making it necessary to shuffle data in registers between LDSM and STSM steps to correctly transpose tiles; details are omitted here.} These instructions are frugal in register usage, enabling their execution by the producer warpgroup, and also support layout transposition during memory movement. Additionally, after the first step, we can schedule the transposition of the following $\mathbf{V}$ tile to be performed concurrently with the two WGMMAs that process the previous $\mathbf{V}$ and current $\mathbf{K}$ tile.

Secondly, it is evident that, in contrast to FP16, the FP32 accumulator employed in an FP8 WGMMA possesses a distinct memory layout from that used for operand A when stored in registers. Portions of these respective layouts are illustrated in \cref{fig:rmem_accum} and \cref{fig:rmem_operand}, where register entries per thread are presented in their corresponding sequences. Utilizing byte permutation instructions allows us to convert the accumulator output from the first WGMMA into a configuration appropriate for the second WGMMA, while aligning with the arrangement of the $\mathbf{V}$ tile generated by the in-kernel transpose. Explicitly, considering \cref{fig:rmem_accum}, we reorder sequentially as follows:
$$\{ \verb|d0 d1 d4 d5 d2 d3 d6 d7| \},$$
and apply this register reshuffling to every 8-byte block. In terms of the abstract structure of the $\mathbf{P}$ tile, this operation rearranges its columns (for instance, columns $0189$ become the leading four columns). To facilitate accurate computation by WGMMA, the in-kernel transpose can be configured to produce a corresponding row rearrangement in the $\mathbf{V}$ tile.\footnote{Having the capability to perform this transpose within the kernel obviates the necessity for shuffle instructions that reallocate register ownership between threads, as discussed previously in~\citep{colfax_fp8_flashattention_2024}.}

\textbf{Accuracy: block quantization and incoherent processing.} The FP8 (e4m3) data type utilizes just 3 bits for the mantissa alongside 4 bits dedicated to the exponent, which inherently leads to greater numerical imprecision compared to FP16/BF16 formats. Additionally, large-scale models are known to exhibit extreme outlier values~\citep{dettmers2208llm,
  sun2024massive}, whose magnitudes significantly exceed those of the majority, complicating the quantization process. A common approach involves per-tensor scaling~\citep{micikevicius2022fp8}, where a distinct scalar is assigned to each tensor (such as a separate value for $\mathbf{Q}$, $\mathbf{K}$, and $\mathbf{V}$).
In our work, to mitigate the inaccuracy associated with attention computations in FP8, we introduce two methods:

\begin{enumerate}

\item \textbf{Block quantization}: This method retains a single scalar for each block; for each of $\mathbf{Q}$, $\mathbf{K}$, and $\mathbf{V}$, the tensor is divided into blocks of size $B_r \times d$ or $B_c \times d$ and each block is quantized independently. Quantization can be efficiently combined with a preceding operation such as rotary embedding, without incurring additional latency, since rotary embedding is limited by memory bandwidth. Because the \textsc{FlashAttention-3} algorithm operates intrinsically on blocks, it is possible to scale each block of $\mathbf{S}$ to compensate for block quantization without additional computation overhead.

\item \textbf{Incoherent processing}: To mitigate the influence of outlier values, both $\mathbf{Q}$ and $\mathbf{K}$ are premultiplied by a random orthogonal matrix $\mathbf{M}$ prior to FP8 quantization. Due to the property that $\mathbf{M}$ is orthogonal ($\mathbf{M} \mathbf{M}^\top = I$), the equality $(\mathbf{Q} \mathbf{M}) (\mathbf{K} \mathbf{M})^\top = \mathbf{Q} \mathbf{K}^\top$ holds, thus the attention mechanism remains unchanged by this transformation. This approach redistributes outliers, as each entry of $\mathbf{Q} \mathbf{M}$ or $\mathbf{K} \mathbf{M}$ corresponds to a randomized linear combination of the original entries, which consequently decreases quantization errors. Consistent with prior work by \citet{chee2024quip} and~\citet{tseng2024quip}, we construct $\mathbf{M}$ as the product of random diagonal matrices with elements $\pm 1$ and a Hadamard matrix, allowing multiplication to be performed in $O(d \log d)$ time instead of $O(d^2)$. Additionally, this multiplication can be merged with the rotary embedding step, incurring no extra computational overhead.
\end{enumerate}

Empirical results indicate that these approaches reduce numerical errors by up to 2.6$\times$, as detailed in \cref{sec:numerical_error}.

\section{Empirical Validation}
\label{sec:exp}

To implement \textsc{FlashAttention-3} and assess its performance and precision, we utilize CUTLASS~\citep{Thakkar_CUTLASS_2023} primitives including WGMMA and TMA abstractions.

\begin{itemize}

\item \textbf{Benchmarking attention.} We evaluate the performance of \textsc{FlashAttention-3} by recording its runtime over various sequence lengths and directly comparing it to several baselines: the conventional PyTorch implementation, \textsc{FlashAttention-2}, the Triton-based version of \textsc{FlashAttention-2} tailored for H100-specific instructions, and the H100 GPU-optimized \textsc{FlashAttention-2} from cuDNN. Our findings demonstrate that \textsc{FlashAttention-3} achieves speedups of up to 2.0$\times$ over \textsc{FlashAttention-2}, and 1.5$\times$ relative to Triton-based \textsc{FlashAttention-2}. Additionally, \textsc{FlashAttention-3} is capable of sustaining throughput up to 740 TFLOPs/s, which corresponds to 75\% of the H100 GPU’s theoretical peak performance.

\item \textbf{Ablation study.} Our experiments indicate that algorithmic enhancements—specifically, warp-specialization and the integration of GEMM-softmax pipelining—are significant factors driving the performance gains observed in \textsc{FlashAttention-3}.

\item \textbf{Accuracy of FP8 attention.} Through empirical validation, we observe that the use of block quantization in conjunction with incoherent processing effectively lowers the numerical error in FP8 \textsc{FlashAttention-3} computations by a factor of 2.6$\times$.

\subsection{Benchmarking Attention}
\label{subsec:benchmark_attn}

We evaluate the execution time of various attention mechanisms using an H100 80GB SXM5 GPU across several configurations, specifically comparing cases with and without causal masking, and head dimensions set to either 64 or 128, all using FP16 input precision. Our findings are presented in~\cref{fig:benchmark_attn_fwd} and~\cref{fig:benchmark_attn_bwd}. The data demonstrate that \textsc{FlashAttention-3} achieves a speedup of roughly 1.5-2.0$\times$ over \textsc{FlashAttention-2} during the forward computation, and is approximately 1.5-1.75$\times$ faster on the backward computation. When placed alongside the typical attention implementation, \textsc{FlashAttention-3} offers accelerations in the range of 3-16$\times$. Moreover, for sequences of intermediate and greater length (1k tokens and above), \textsc{FlashAttention-3} even outperforms proprietary libraries like cuDNN—designed and tuned for the H100 architecture.

\paragraph{Benchmark settings:} We vary the sequence length as 512, 1k, ..., 16k, and set batch size so that the total number of tokens is 16k. We set the hidden dimension to 2048, and head dimension to be either 64, 128, or 256 (i.e., 32 heads, 16 heads, or 8 heads). To calculate the FLOPs of the forward pass, we use:

\begin{equation*}
  4 \cdot \text{seqlen}^2 \cdot \text{head dimension} \cdot \text{number of heads}.
\end{equation*}

Applying causal masking involves halving the calculated value, reflecting that roughly 50% of the matrix elements are actually processed. For the backward pass, the computational load in terms of FLOPs is estimated by scaling the forward pass FLOPs by a factor of 2.5; this ratio emerges from the presence of 2 matrix multiplications during the forward pass and 5 matrix multiplications required in the backward pass, including those due to recomputation.

Additionally, we assess the forward pass runtime for FP8 using analogous configurations. The outcomes for headdim 256 are presented in ~\cref{fig:benchmark_attn_fp8}, while comprehensive results can be found in \cref{sec:benchmark_attn_fp8_full}.

\subsection{Ablation Study: 2-Stage Pipelining Experiments}
\label{sec:2-stage-pipelining-experiments}

We conduct ablation experiments on both the 2-stage WGMMA-softmax pipelining and warp-specialization in the context of non-causal FP16 \textsc{FlashAttention-3}, maintaining fixed parameter settings of $\{\text{batch}, \text{seqlen}, \text{nheads}, \text{hdim}\} = \{ 4, 8448, 16, 128\}$. As shown in~\cref{table:ablation_pipelining}, the outcomes validate that our algorithmic enhancements—specifically, asynchronous execution leveraging warp-specialization and overlapping the GEMM and softmax computations—substantially boost performance, increasing throughput from 570 to 661 TFLOPs.

\subsection{Numerical Error Validation}
\label{sec:numerical_error}

Given the ongoing concern regarding numerical error~\citep{golden2024flash} in
\textsc{FlashAttention}, we examine \textsc{FlashAttention-2}, \textsc{FlashAttention-3}, and a conventional attention implementation, comparing each to a reference version using FP64 precision. To mimic atypical features and activations found in LLMs~\citep{dettmers2208llm, sun2024massive}, we create the elements of $\mathbf{Q}, \mathbf{K}, \mathbf{V}$ according to the subsequent distribution:

\begin{equation*}
  \mathcal{N}(0, 1) + \mathcal{N}(0, 100) \cdot \mathrm{Bernoulli}(0.001).
\end{equation*}

Specifically, each element is sampled from a normal distribution with mean zero and unit standard deviation, except for 0.1\% of entries where an additional independent Gaussian term with standard deviation 10 is introduced. The root mean squared error (RMSE) obtained under these conditions is reported in~\cref{table:numerical_error}. When evaluated in FP16, both \textsc{FlashAttention-2} and \textsc{FlashAttention-3} achieve a reduction in RMSE by a factor of 1.7$\times$ relative to the default implementation, attributable to the use of FP32 precision for intermediate (softmax) computations. The baseline FP8 attention utilizes per-tensor scaling; its matrix multiplication accumulator operates in FP32, while softmax intermediates are maintained in FP16. By leveraging block quantization and incoherent computation, \textsc{FlashAttention-3} in FP8 demonstrates 2.6$\times$ greater accuracy than the baseline.

\section{Dicussion, Limitations, Conclusion}
\label{sec:discussion}

By leveraging \textsc{FlashAttention-3}, we showcase how advancements in programming paradigms and hardware capabilities—including asynchronous execution and reduced-precision formats—can significantly enhance both the performance and accuracy of attention mechanisms. Our implementation achieves a 1.5-2.0$\times$ increase in attention speed over \textsc{FlashAttention-2}, and exhibits a 2.6$\times$ improvement in FP8 numerical error when compared with conventional per-tensor quantization strategies. Several aspects of this work remain open for further refinement: enhancing optimization for LLM inference scenarios, incorporating a persistent kernel approach within the FP8 kernel,\footnote{In our experiments, FP16 \textsc{FlashAttention-3} utilizes both a persistent kernel and a load balancing mechanism, whereas FP8 \textsc{FlashAttention-3} currently does not. This contributes in part to the comparatively lower performance of FP8 \textsc{FlashAttention-3} for small sequence lengths and causal masking, relative to FP8 cuDNN kernels.} and exploring the influence of reduced precision in large-scale training settings. Although this study centers on Hopper GPUs, we anticipate that the methodologies outlined here may be applicable to a variety of hardware accelerators. Ultimately, we aspire that the advancement of faster and more reliable primitives like attention could enable emerging possibilities for tasks requiring extended context.

\end{document}